\section{Grundlagen}
Dieses Kapitel beschreibt die grundlegenden technischen Rahmenbedingungen der umgesetzten Lösung. Dabei werden die Systemarchitektur sowie die verwendeten Technologien als Basis für die nachfolgenden Implementierungs- und Funktionsbeschreibungen vorgestellt.

\subsection{Systemarchitektur}
\subsection{Technologien}

Dieses Unterkapitel gibt einen Überblick über die zentralen Technologien, die für die Umsetzung des Systems verwendet wurden. Dabei werden insbesondere Backend, Datenhaltung, Suchtechnologie sowie Frontend-nahe Technologien kurz eingeordnet.

\subsection*{Backend und Kommunikation}
Die Backend-Implementierung erfolgte in der Programmiersprache \textbf{C\#}. Zur Kommunikation zwischen Client und Server wird eine \textbf{REST-basierte Schnittstelle (REST-API)} verwendet. Diese ermöglicht eine standardisierte Datenübertragung über \textbf{HTTP} und sorgt für eine klare Trennung zwischen Präsentationsschicht und Anwendungslogik.

\subsection*{Datenhaltung}
Zur persistenten Speicherung der Daten kommt das relationale Datenbankmanagementsystem \textbf{PostgreSQL} zum Einsatz. PostgreSQL bietet eine hohe Stabilität, umfangreiche SQL-Funktionalitäten sowie eine gute Skalierbarkeit und eignet sich daher optimal für datenbankgestützte Webanwendungen.

\subsection*{Suchtechnologie}
Für die Dokumentensuche wird \textbf{Azure AI Search} eingesetzt. Dieser Dienst ermöglicht die Indizierung großer Dokumentbestände sowie eine performante Volltextsuche mit Filter- und Facettenfunktionen. Dadurch können Dokumente anhand von Suchbegriffen und strukturierten Metadaten effizient aufgefunden werden.

\subsection*{Frontend und Design}
Das Frontend wurde mit einem \textbf{Blazor WebAssembly} realisiert. Dieses Framework ermöglicht die Entwicklung interaktiver Webanwendungen unter Verwendung von C\#, wobei der Anwendungscode direkt im Browser des Clients ausgeführt wird. Dadurch kann ein Teil der Verarbeitungslogik clientseitig erfolgen, was zu einer verbesserten Benutzerinteraktion führt.

\begin{itemize}
    \item \textbf{CSS:} Das Layout und Design der Benutzeroberfläche wird mithilfe von Cascading Style Sheets (CSS) umgesetzt. Dadurch können Struktur, Darstellung und Responsivität flexibel gestaltet werden.
    \item \textbf{JavaScript:} Zusätzlich wird JavaScript für spezifische browserseitige Funktionen verwendet. Insbesondere kommt es beim Upload von Dokumenten zum Einsatz, um browsernahe Funktionen effizient zu nutzen und die Interaktion mit Dateisystemen zu ermöglichen.
\end{itemize}
